\documentclass[../main.tex]{subfiles}
\begin{document}
\section{Introducción a la Criptografía}

Esta unidad usa la \textbf{historia de la criptografía} como vehículo para
introducir, de forma intuitiva, los conceptos fundamentales: criptosistema,
clave, principio de Kerckhoffs, criptoanálisis (ataque por fuerza bruta y
análisis estadístico) y la noción de seguridad. Cierra con una primera
formalización (terna de algoritmos) y la definición de \textbf{secreto
perfecto} de Shannon, cuyo desarrollo formal continúa en la unidad 2.

\prof{Hoy vamos a ver fundamentos, vamos a ver ideas. No vamos a entrar en el campo formal. Vamos a aprovechar la historia de la criptografía como para ir entendiendo los conceptos básicos.}

\begin{examen}
Los cifrados históricos son el \textbf{vehículo recurrente} para discutir
secreto perfecto. En distintos años se pidió analizar: Hill (2015),
transposición por columnas (2016), cilindro de Jefferson (2017), homofónico
(2018), Vigenère (2025-2). Tareas típicas: \textbf{calcular los tamaños de los
espacios} $\#M$, $\#K$, $\#C$ y \textbf{demostrar o refutar secreto perfecto}.
También caen: principio de Kerckhoffs (V/F), análisis de frecuencias e índice
de coincidencia, y distinguir sustitución vs.\ transposición / confusión vs.\
difusión / linealidad.
\end{examen}

%======================================================================
\subsection{Qué es la criptografía}

Del griego: \textbf{escritura secreta}. Nace de la necesidad de ocultar
información (``desde que el hombre es hombre y se agrupa, existe la necesidad de
mantener alguna cosa en secreto'').

\begin{definicion}{Criptografía}
Conjunto de \textbf{funciones matemáticas y técnicas} que permiten asegurar
distintas propiedades sobre la información. Son las herramientas básicas desde
donde se construye la seguridad.
\end{definicion}

El concepto moderno trasciende la sola \textbf{confidencialidad} (esconder
información) y abarca otras propiedades. Usos cotidianos donde hay criptografía
``oculta en todos lados'': tráfico web e inalámbrico (Wi-Fi, satélite),
protección de archivos en disco, autenticación de usuarios (segundo factor),
protección de contenido (DRM, e-readers), firmas digitales, voto electrónico y
dinero electrónico (Bitcoin / criptomonedas).

\medskip
Existe una distinción de especialidades:
\begin{itemize}[nosep]
  \item \textbf{Criptografía}: estudio (diseño) de las funciones.
  \item \textbf{Criptoanálisis}: estudio de los \textbf{ataques} a esas funciones.
\end{itemize}

\begin{destacado}
Objetivo secundario de la materia: entender la \textbf{complejidad} que hay
detrás de las funciones criptográficas para \textbf{no inventar primitivas
propias}. El profe estima que hay unas 30 o 40 personas en el mundo capaces de
construir funciones que sobrevivan al escrutinio público. Como ingenieros,
debemos \textbf{usar las que ya existen}.
\end{destacado}

%======================================================================
\subsection{Criptosistema (visión informal)}

La construcción más básica y común es el \textbf{criptosistema}: en principio,
un conjunto de dos funciones.

\begin{itemize}[nosep]
  \item \textbf{Cifrado / encriptar}: toma un mensaje (texto plano) y lo
        transforma en algo aparentemente sin sentido (texto cifrado).
  \item \textbf{Descifrado / desencriptar}: toma el mensaje cifrado y recupera
        el mensaje original.
\end{itemize}

\paragraph{Función de cifrado.} Es una función de \textbf{dos parámetros}:
$$\text{Mensaje plano} \times \text{clave} \to \text{Mensaje cifrado}, \qquad e_k(p) = c.$$
Diferentes representaciones, mismo significado:
$$e_k(p) = \mathrm{enc}_k(p) = e(k,p) = \{p\}_k.$$
(La notación $\{p\}_k$ se usa más en protocolos: lo que está entre llaves está
cifrado con la clave del subíndice.)

\paragraph{Función de descifrado.} También de dos parámetros:
$$\text{Mensaje cifrado} \times \text{clave} \to \text{Mensaje plano}, \qquad d_k(c) = p.$$

\paragraph{La clave.} Es un \textbf{bloque arbitrario de información}. Se la
separa del mensaje y se le da un nombre especial porque, al probar la seguridad
de un sistema, \textbf{siempre se presupone que la clave está protegida}.

\begin{destacado}
Ejercicio mental fundamental en seguridad: \textbf{ponerse en el lugar del
atacante}. La única constante es que, donde aparezca algo llamado \emph{clave},
se asume que el atacante \textbf{no la tiene}; y si la consigue, ``se acabó la
seguridad''.
\end{destacado}

\paragraph{¿Para qué sirve?} Resuelve el problema de \textbf{controlar el acceso
y la diseminación} de la información. La idea: ``patear el problema''. Proteger
un mensaje de tamaño arbitrario es difícil; en cambio, ciframos el mensaje con
una clave y entonces el texto cifrado puede almacenarse o transmitirse por
cualquier lado (es público), mientras solo haya que \textbf{proteger la clave},
que es mucho más chica. En los criptosistemas modernos la clave \textbf{no
depende del tamaño del mensaje}: se pueden cifrar terabytes con una clave de
128 bits.

%======================================================================
\subsection{Principio de Kerckhoffs}

\begin{definicion}{Principio de Kerckhoffs}
Un criptosistema debe ser \textbf{seguro incluso si todo sobre el sistema}
(todo el algoritmo, todas las funciones), \textbf{excepto la clave}, es de
público conocimiento.
\end{definicion}

Consecuencia para el análisis: la seguridad \textbf{no puede depender} de que
el atacante desconozca el algoritmo o el código fuente. Hay que asumir que el
atacante \textbf{conoce el criptosistema usado} y eventualmente el código; lo
único secreto es la clave.

\begin{ojo}
Error clásico de V/F. La afirmación ``la seguridad solo depende de que el
\emph{algoritmo} esté protegido'' es \textbf{FALSA} (Recup.\ 2025). Lo correcto:
la seguridad depende \textbf{solo de la clave}, no de mantener secreto el
algoritmo.
\end{ojo}

%======================================================================
\subsection{Seguridad (noción informal)}

Primera intuición de un alumno: ``es seguro si quien no tiene la clave no puede
descifrar el mensaje en tiempo razonable''. El profe la refina: se pide
\textbf{más}.

\begin{definicion}{Seguridad (informal)}
Un criptosistema es seguro si \textbf{ningún adversario puede computar
\emph{cualquier} función del texto plano} a partir del mensaje cifrado que posee.
Significa que no se puede recuperar:
\begin{itemize}[nosep]
  \item el mensaje,
  \item parte del mensaje,
  \item ni siquiera \textbf{el sentido} del mensaje.
\end{itemize}
\end{definicion}

\prof{No es solamente que no lo puedas decodificar: es mucho más fuerte, no se puede extraer ningún tipo de información a partir del texto cifrado si no se tiene la clave.}

Ejemplos de por qué importa ``parte'' y ``sentido'': un documento gigante cuyo
resumen está en la primera página; o una negociación secreta donde al adversario
le basta saber si la respuesta será \emph{positiva} o \emph{negativa} (sin
descifrar todo) para obtener una ventaja ilegal.

%======================================================================
\subsection{Breve historia de la criptografía}

\begin{center}
\begin{tabular}{ll}
\textbf{$\sim$2500 AC} & Egipto: sustitución (jeroglíficos no estándar). \\
\textbf{$\sim$700--300 AC} & Escítala (Scitala): transposición con báculos. \\
\textbf{50 AC} & César: rotación (subtipo de sustitución). \\
\textbf{800 DC} & Primeros documentos de criptoanálisis: análisis de frecuencias. \\
\textbf{$\sim$1500} & Sustitución \textbf{polialfabética} (un símbolo cifrado \\
                    & representa distintos símbolos del original). \\
\textbf{1553} & Cifrado de \textbf{Vigenère}. \\
\textbf{1939} & Enigma / Bombe: ataques de fuerza bruta (protocomputadoras). \\
\textbf{1949} & Shannon: \emph{Information Theory \& Cryptography}, \\
              & inicio de la \textbf{criptografía moderna}. \\
\end{tabular}
\end{center}

\begin{itemize}[nosep]
  \item \textbf{Escítala}: báculo de cierto tamaño/topología; se enrolla una
        tira de cuero y se escribe transversalmente. Para descifrar se necesita
        el \textbf{mismo báculo}. Es \textbf{transposición}.
  \item \textbf{César}: rotación \textbf{fija de 3} posiciones en el alfabeto
        romano. Documentado en campañas militares.
  \item Las dos grandes revoluciones: (1) los polialfabéticos
        ($\sim$400 años de dominio); (2) la criptografía moderna, contemporánea
        a la 2.ª Guerra Mundial (Shannon, teoría de la información).
\end{itemize}

%======================================================================
\subsection{Cifrado por rotación (ROT-X / César)}

Se asigna un orden a los símbolos (alfabeto, comenzando por 0). La clave $k$ es
un número entre 1 y 26 (cantidad de letras $-1$). Cifrar = reemplazar cada letra
por la que ocupa $k$ posiciones más adelante, de forma \textbf{cíclica} (después
de la \texttt{z} vuelve a la \texttt{a}).

$$e(\text{prueba}, 4) = \text{tvyife}$$

\paragraph{Criptoanálisis.} \textbf{Debilidad: hay pocas claves.} Conociendo el
algoritmo (Kerckhoffs) se prueban todas las claves por \textbf{fuerza bruta}:
\begin{center}
\begin{tabular}{lll}
$k=1$ & $d(1,\text{tvyife})=\text{suxhed}$ & sin sentido \\
$k=2$ & $d(2,\text{tvyife})=\text{rtwgdc}$ & sin sentido \\
$k=3$ & $d(3,\text{tvyife})=\text{qsvfcb}$ & sin sentido \\
$k=4$ & $d(4,\text{tvyife})=\text{prueba}$ & ¡con sentido! \\
\end{tabular}
\end{center}

\begin{destacado}
El ataque por \textbf{fuerza bruta siempre está disponible} para el atacante,
incluso contra el criptosistema más complejo: dado el texto cifrado y conocido
el algoritmo, nada impide probar todas las claves. La debilidad de ROT-X no es
que se pueda hacer fuerza bruta, sino que las \textbf{claves son pocas}.
\end{destacado}

\begin{ojo}
La fuerza bruta tiene un \textbf{prerrequisito}: poder \textbf{discriminar un
descifrado válido de uno inválido}. Si el mensaje es una sola letra elegida con
distribución uniforme, el cifrado puede ser cualquier letra y \textbf{no se
puede saber} cuál era el original (los descifrados son todos igual de
``válidos''). Este es el germen de la idea de secreto perfecto.
\end{ojo}

\paragraph{Redundancia y distancia de unicidad (Shannon).} Los lenguajes
naturales tienen mucha \textbf{redundancia} (mecanismo de control de errores:
podemos leer una frase sin vocales). La \textbf{redundancia} es la diferencia
entre el tamaño del mensaje y la información que realmente lleva (comprimir =
recodificar eliminando redundancia; p.\ ej.\ Huffman). La \textbf{distancia de
unicidad} es la longitud mínima que debe tener un mensaje para que la
probabilidad de que exista más de un descifrado con sentido sea despreciable.
Para lenguajes naturales (español, inglés) está en el orden de los
\textbf{4 caracteres}: por encima de ese umbral, la fuerza bruta que encuentra
un texto con sentido suele dar el resultado correcto.

%======================================================================
\subsection{Cifrado por sustitución (monoalfabético)}

Para cada símbolo del alfabeto se establece una \textbf{transformación única}
(una permutación del alfabeto): no puede haber dos símbolos distintos que se
transformen en lo mismo (si no, no sería invertible). Reemplaza un símbolo por
otro (que puede ser de otro alfabeto: letras, símbolos, números).

\begin{center}\ttfamily
\begin{tabular}{l}
k = dublcmfthijnzpxqeaosvkrwgy \\
\ \ \ \ abcdefghijklmnopqrstuvwxyz \\[4pt]
esto es una prueba \;$\to$\; cosx co vpd qavcud \\
\end{tabular}
\end{center}

\paragraph{Tamaño del espacio de claves.} La clave es la permutación del
alfabeto. Con el alfabeto español de \textbf{27 símbolos}:
$$\#K = 27! \approx \text{número de 28 dígitos (cuatrillones).}$$
Es un número ``bestialmente alto'': la fuerza bruta es inviable. \textbf{No}
tiene la debilidad de ROT-X (pocas claves).

\begin{examen}
``¿Cuántas claves tiene este criptosistema?'' Para sustitución monoalfabética
sobre un alfabeto de $n$ símbolos: $\#K = n!$ (es una \textbf{permutación}).
Para el español, $n=27 \Rightarrow 27!$.
\end{examen}

\paragraph{Criptoanálisis: análisis de frecuencias.}
\begin{ojo}
\textbf{Debilidad distinta}: la sustitución \textbf{preserva las características
estadísticas del lenguaje}. Una misma letra del plano siempre va al mismo
símbolo cifrado, así que la letra que más aparece en el plano es la que más
aparece en el cifrado.
\end{ojo}
Procedimiento (resoluble incluso por un chico, o sin conocer el idioma usando
solo los indicadores estadísticos):
\begin{enumerate}[nosep]
  \item Obtener la \textbf{tabla de frecuencias} estimadas del lenguaje.
  \item Calcular la frecuencia de cada símbolo en el texto cifrado.
  \item Asumir que los símbolos más probables se corresponden (en español: E,
        A, S, \dots).
  \item Usar grupos de 2 y 3 letras comunes (\emph{el}, \emph{la}, \emph{de},
        \emph{los}, \emph{que}\dots), reglas (\emph{q} siempre seguida de
        \emph{u}; no hay 3 consonantes juntas), reemplazar y verificar sentido.
\end{enumerate}

\paragraph{Tabla de frecuencias del español (\%).}
\begin{center}
\begin{tabular}{lll}
\multicolumn{3}{l}{\textbf{Frecuencia alta}} \\
E 13,11 & A 10,60 & S 8,47 \\
O 8,23 & I 7,16 & N 7,14 \\
R 6,95 & D 5,87 & T 5,40 \\[4pt]
\multicolumn{3}{l}{\textbf{Frecuencia media}} \\
C 4,85 & L 4,42 & U 4,34 \\
M 3,11 & P 2,71 & G 1,40 \\
B 1,16 & F 1,13 & V 0,82 \\[4pt]
\multicolumn{3}{l}{\textbf{Frecuencia baja}} \\
Y 0,79 & Q 0,74 & H 0,60 \\
Z 0,26 & J 0,25 & X 0,15 \\
W 0,12 & K 0,11 & Ñ 0,10 \\
\end{tabular}
\end{center}

\begin{destacado}
La debilidad estadística \textbf{se traslada al mundo digital}: cifrar un MP4,
un Word, etc., también deja estructura estadística (e incluso \emph{headers} con
valores fijos), atacable con los mismos mecanismos. Una cosa es decir que un
criptosistema es \textbf{débil / fácil de romper} y otra es \textbf{efectivamente
romperlo}: cuanto más complejo el sistema, más grande el espacio de búsqueda.
\end{destacado}

\begin{ojo}
``Es seguro porque nadie lo rompió'' es un criterio \textbf{malo}. Muchos
criptosistemas fueron considerados buenos solo porque nadie les había dedicado
el tiempo necesario, no por ser fundacionalmente seguros.
\end{ojo}

\paragraph{Gancho (crib / known-plaintext).} Saber que en el plano aparece una
secuencia conocida (p.\ ej.\ un saludo fijo al inicio, o ``LACABEZA'') acelera
mucho el ataque. Históricamente, los aliados aprovecharon que los mensajes
alemanes empezaban casi siempre con el mismo saludo.

%======================================================================
\subsection{Cifrado de Vigenère (sustitución polialfabética)}

Es una \textbf{composición de cifrados de rotación}: un mismo símbolo del plano
puede transformarse en \textbf{distintos} símbolos cifrados. La clave es una
secuencia de $n$ números (o letras, con A$=0$, B$=1$, \dots); cada elemento
indica cuántas posiciones rotar la letra correspondiente. Al terminar la clave,
se vuelve a empezar (el mensaje se divide en \textbf{bloques} del largo de la
clave).

\begin{center}\ttfamily
\begin{tabular}{l}
K = ECFD = 4253 \quad (long 4) \\
esto es una prueba \\
4253 42 534 253425 342 \dots \\
$\to$ iuyrdgxcyofcttzhfc \\
\end{tabular}
\end{center}

\begin{destacado}
Vigenère \textbf{rompe las propiedades estadísticas} de la sustitución simple:
la misma letra del plano puede dar letras cifradas distintas, y el mismo símbolo
cifrado puede provenir de letras distintas. La letra que más aparece deja de ser
el símbolo más frecuente del cifrado. Creado en \textbf{1553}, se lo consideró
seguro casi \textbf{300 años}, hasta que en \textbf{1863 Friedrich Kasiski}
publicó un método para resolverlo.
\end{destacado}

\paragraph{Criptoanálisis: Test de Kasiski (dos fases).}

\textbf{Fase 1 — determinar la longitud de la clave (= tamaño de bloque).}
Aparecen \textbf{n-gramas repetidos} cuando los mismos símbolos del plano caen
en la misma posición dentro del bloque (se rotan igual). Entonces:
\begin{enumerate}[nosep]
  \item Buscar secuencias repetidas en el texto cifrado.
  \item Calcular la \textbf{distancia} (en caracteres) entre cada par de
        apariciones.
  \item Esas distancias deben ser \textbf{múltiplos} del tamaño de bloque, así
        que la longitud de la clave es el \textbf{MCD} de las distancias halladas.
\end{enumerate}
Ejemplo: GGMP (×3) separadas 256 y 104; YEDS (×2) separada 72; HASE (×2)
separada 156; VSUE (×2) separada 32.
$$\mathrm{MCD}(32, 72, 104, 156, 256) = 4.$$
Es una prueba estadística: puede fallar si dos secuencias coinciden por
casualidad (el MCD da 1); en ese caso se van descartando muestras de a una hasta
que el MCD vuelva a subir.

\textbf{Fase 2 — resolver cada bloque (método de coincidencia mutua).}
Conocida la longitud (p.\ ej.\ 4), se agrupan las letras por posición:
todas las 1.ª de cada bloque, todas las 2.ª, etc. Cada grupo está rotado por
\textbf{una única} cantidad, así que su histograma de frecuencias se parece al
del lenguaje. Idea de Kasiski:
\begin{itemize}[nosep]
  \item Al \textbf{juntar} dos grupos con rotaciones distintas, el histograma se
        rompe.
  \item Pero existe una rotación relativa que los \textbf{realinea}: rotando un
        grupo respecto del otro hasta que el histograma combinado se mantenga,
        se obtiene la \textbf{diferencia relativa de rotación} entre posiciones.
  \item Encadenando todos los grupos, el problema se reduce a un \textbf{único
        cifrado por rotación} (todo el mensaje desplazado por el mismo número),
        que se resuelve probando las $\sim$26 posibilidades.
\end{itemize}

\begin{destacado}
A partir de Vigenère, los ataques dejan de ser genéricos: \textbf{ya no hay una
solución única} para toda la familia polialfabética. El método de Kasiski
\textbf{explota el conocer el algoritmo de Vigenère}; cada criptosistema
polialfabético requiere un ataque \emph{ad hoc}. (Para sustitución simple sí
existía un algoritmo genérico.)
\end{destacado}

%======================================================================
\subsection{Máquina Enigma (rotores)}

Estado del arte en la 2.ª Guerra Mundial: criptosistema polialfabético de
\textbf{rotores} usado por Alemania. Cada rotor es como una sustitución simple
(entra una letra, sale otra); al rotar, el mapeo cambia. Enigma original tenía
\textbf{3 rotores} (hubo versiones de 4):
\begin{itemize}[nosep]
  \item Cada vez que se tipea una letra, rota un rotor; al dar la vuelta, rota el
        siguiente (como un cuentakilómetros).
  \item Período de repetición $\approx 26^3$, mucho mayor que cualquier mensaje:
        en la práctica, la secuencia \textbf{no se repetía}. Equivale a una
        sustitución simple que \textbf{varía letra a letra}, con ``bloque'' de
        $26^3$ en vez de los $4$ de Vigenère.
  \item Un \textbf{plugboard} (tablero) cruzaba conexiones; un reflector hacía
        que la señal ``rebotara'', de modo que la \textbf{misma máquina} servía
        para cifrar y descifrar (poniendo la misma clave inicial).
  \item La \textbf{clave} era la disposición inicial de los rotores (cambiaba
        cada día según un cuaderno por oficial).
\end{itemize}
Los aliados (liderados en el criptoanálisis por \textbf{Turing}, entre otros)
hicieron \textbf{ingeniería inversa} tras capturar una máquina y construyeron
una computadora (Bombe) para fuerza bruta. No alcanzaba el cómputo; lo
resolvieron usando \textbf{ganchos} (mensajes que empezaban con un saludo
conocido) para reducir el espacio de búsqueda. \textbf{No rompieron Enigma de
forma general}, solo ese uso; los submarinos con 4 rotores no fueron descifrados
durante la guerra. De este esfuerzo surge prácticamente la computadora moderna
(arquitectura de von Neumann).

%======================================================================
\subsection{Hacia la criptografía moderna}

La 2.ª Guerra generó una \textbf{crisis}: no había criterio para dar fiabilidad
de la seguridad. La definición era ``es seguro porque nadie lo rompió'', y al
juntar gente brillante se rompieron cosas que se creían seguras. De ahí nacen
los \textbf{tres pilares} de la criptografía moderna:
\begin{enumerate}[nosep]
  \item \textbf{Representaciones formales} de los criptosistemas (ya no recetas
        algorítmicas).
  \item Definiciones del \textbf{modelo de amenaza} (no hay una única definición
        de ``seguro''; Enigma quizá era seguro, pero \emph{como lo usaban} no lo
        era).
  \item \textbf{Pruebas formales} de seguridad.
\end{enumerate}
Los cifrados clásicos vistos \textbf{no se usan en la práctica} (son fáciles de
romper algorítmicamente), pero son la \textbf{base} sobre la que se construye lo
moderno.

\paragraph{Criptosistema (primera definición formal).}
\begin{definicion}{Criptosistema (terna de algoritmos)}
\begin{itemize}[nosep]
  \item $\text{Gen}: () \to K$ \quad — \textbf{generador de clave}. Suele ser
        sin parámetros; \emph{determina el espacio de claves} (cuántas y cuáles
        son posibles).
  \item $\text{Enc}: K \times P \to C$ \quad — \textbf{cifrado}.
  \item $\text{Dec}: K \times C \to P$ \quad — \textbf{descifrado}.
\end{itemize}
\textbf{Propiedad básica (corrección):} para todo $m$ y $k$ válidos,
$$d_k(e_k(m)) = m.$$
\end{definicion}
Conjuntos: $K$ (espacio de claves), $P$ o $M$ (espacio plano, todos los
mensajes), $C$ (espacio cifrado).

\begin{ojo}
La propiedad básica es solo ``sanitaria'': la \textbf{función identidad}
(ignorar la clave y devolver el mensaje) también la cumple. Por eso esta
definición \textbf{no} habla todavía de seguridad — porque no hay una única
noción de seguridad.
\end{ojo}

%======================================================================
\subsection{Secreto perfecto (introducción)}

Shannon define el primer criterio formal de seguridad: el \textbf{secreto
perfecto}.

\begin{definicion}{Secreto perfecto (Shannon)}
Dado un criptosistema $(\text{Gen}, e, d)$, posee secreto perfecto si para
\textbf{toda} distribución de probabilidades en $M$, cada mensaje $m$ y cada
texto cifrado $c$ con $\Pr[C=c] > 0$:
$$\Pr[M=m \mid C=c] = \Pr[M=m].$$
\end{definicion}

Interpretación: la probabilidad que se conoce \emph{a priori} sobre el mensaje
\textbf{no se modifica} al conocer el texto cifrado. Conocer $c$ no aporta
\textbf{ninguna información} sobre $m$.

\prof{Es mucho más fuerte que no poder decodificar el mensaje: no podés extraer nada de información. Un criptosistema con secreto perfecto es imposible de romper, incluso por fuerza bruta.}

Ejemplo (mensaje sí/no con $\Pr[\text{sí}]=0{,}8$ a priori): si el sistema tiene
secreto perfecto, decirte el cifrado no cambia ese $0{,}8$ a $0{,}81$ ni a
$0{,}79$.

\begin{ojo}
\textbf{Error típico de examen:} confundir ``secreto perfecto'' con ``no se
puede romper / decodificar''. El criterio es la \textbf{definición formal de
Shannon}, $\Pr[M=m\mid C=c]=\Pr[M=m]$, que es \textbf{más fuerte}: no se aprende
\emph{nada}. Como consecuencia (se demuestra en la unidad 2), exige
$$\#K \ge \#M.$$
\end{ojo}

\begin{destacado}
\textbf{Paradoja del secreto perfecto:} como $c = e_k(m)$, las variables
aleatorias $C$ y $M$ están \textbf{relacionadas por una función} (son
dependientes por diseño); sin embargo, la condición de secreto perfecto pide que
se comporten como \textbf{estadísticamente independientes}
($\Pr[M=m\mid C=c]=\Pr[M=m]$). Queda la pregunta: ¿es una asíntota teórica
inalcanzable, o existe una construcción que lo logra? (Se retoma en la unidad 2.)
\end{destacado}

\begin{examen}
Esquema típico de problema con un cifrado histórico (Hill, transposición por
columnas, Jefferson, homofónico, Vigenère):
\begin{enumerate}[nosep]
  \item Calcular $\#M$, $\#K$, $\#C$ (p.\ ej.\ permutaciones $\Rightarrow n!$;
        rotación $\Rightarrow \#K = n-1$; Vigenère de largo $\ell$ sobre $n$
        símbolos $\Rightarrow \#K = n^{\ell}$).
  \item Verificar la condición necesaria $\#K \ge \#M$.
  \item Demostrar o refutar secreto perfecto vía
        $\Pr[M=m\mid C=c]=\Pr[M=m]$ (o exhibiendo $m,c$ que la violan).
\end{enumerate}
\end{examen}

\medskip
\noindent\textit{Lectura recomendada:} Katz \& Lindell, \emph{Introduction to
Modern Cryptography}, Capítulo 1.

\end{document}
